\section{Regulacje i wymagania dla systemów eClinical}

% SLIDE ID: reg-systems-goal
\BlockGoalFrame{\SlideID{reg-systems-goal}}{
  \item zobaczyć, jakie wymagania dotyczą systemów eClinical,
  \item połączyć regulacje z praktycznymi funkcjami systemu,
  \item zrozumieć, że zgodność to proces, nie tylko funkcja w aplikacji.
}

% SLIDE ID: reg-main-sources
\begin{frame}{Główne źródła wymagań}
  \SlideID{reg-main-sources}
  \begin{itemize}
    \item ICH E6(R3) Good Clinical Practice,
    \item EMA guideline on computerised systems and electronic data,
    \item FDA 21 CFR Part 11 i guidance 2024,
    \item EU Clinical Trials Regulation 536/2014,
    \item lokalne SOP, umowy i walidacja systemu.
  \end{itemize}
\end{frame}

% SLIDE ID: reg-system-scope
\begin{frame}{Jakie systemy obejmują wymagania}
  \SlideID{reg-system-scope}
  \begin{itemize}
    \item EDC i eCRF,
    \item IVRS / IWRS / IRT,
    \item eTMF i repozytoria dokumentów,
    \item CTMS i systemy monitoringu,
    \item ePRO, eCOA, aplikacje pacjenta i urządzenia.
  \end{itemize}

  \vspace{0.3cm}

  \begin{ExampleBox}
    Regulator patrzy nie tylko na aplikację, ale na cały cykl życia danych: od wpisu do archiwizacji.
  \end{ExampleBox}
\end{frame}

% SLIDE ID: reg-core-principles
\begin{frame}{Wspólne wymagania}
  \SlideID{reg-core-principles}
  \begin{itemize}
    \item system powinien być zwalidowany proporcjonalnie do ryzyka,
    \item dostęp musi być kontrolowany i przypisany do osoby,
    \item dane muszą być kompletne, czytelne, dostępne i chronione,
    \item zmiany powinny być widoczne w audit trail,
    \item procedury muszą opisywać użycie, backup, awarie i archiwizację.
  \end{itemize}
\end{frame}

% SLIDE ID: reg-data-integrity
\begin{frame}{Integralność danych}
  \SlideID{reg-data-integrity}
  \begin{itemize}
    \item dane muszą być przypisywalne do użytkownika,
    \item wpis powinien być rejestrowany w czasie zdarzenia,
    \item oryginalne dane i metadane muszą pozostać dostępne,
    \item system nie może pozwalać na niewidoczne zmiany „od tyłu”.
  \end{itemize}

  \vspace{0.3cm}

  \begin{QuestionBox}
    Czy screenshot z systemu jest wystarczającym dowodem danych?
  \end{QuestionBox}
\end{frame}

% SLIDE ID: reg-edc-ecrf
\begin{frame}{EDC / eCRF}
  \SlideID{reg-edc-ecrf}
  \begin{itemize}
    \item formularze powinny zbierać dane wymagane protokołem,
    \item walidacje i edit checks wspierają kompletność danych,
    \item query musi zostawiać ślad decyzji i odpowiedzi,
    \item audit trail powinien pokazywać kto, kiedy i co zmienił.
  \end{itemize}

  \vspace{0.3cm}

  \begin{ExampleBox}
    Dobre eCRF nie zbiera „wszystkiego”, tylko dane potrzebne do oceny bezpieczeństwa i punktów końcowych.
  \end{ExampleBox}
\end{frame}

% SLIDE ID: reg-ivrs-iwrs
\begin{frame}{IVRS / IWRS / IRT}
  \SlideID{reg-ivrs-iwrs}
  \begin{itemize}
    \item randomizacja musi być możliwa do odtworzenia,
    \item zaślepienie musi być chronione technicznie i proceduralnie,
    \item dostęp do unblindingu powinien być ograniczony,
    \item przypisanie leku musi być spójne z protokołem i zapasami.
  \end{itemize}

  \vspace{0.3cm}

  \begin{WarningBox}
    Największe ryzyko: błąd systemu lub użytkownika, który narusza randomizację albo zaślepienie.
  \end{WarningBox}
\end{frame}

% SLIDE ID: reg-etmf
\begin{frame}{eTMF}
  \SlideID{reg-etmf}
  \begin{itemize}
    \item TMF musi umożliwiać ocenę przebiegu badania,
    \item dokumenty powinny być kompletne, aktualne i dostępne,
    \item wersje, podpisy i daty muszą być kontrolowane,
    \item archiwizacja musi zachować czytelność i dostępność przez wymagany okres.
  \end{itemize}

  \vspace{0.3cm}

  \begin{QuestionBox}
    Czy eTMF to tylko „dysk na dokumenty”, czy system kontroli badania?
  \end{QuestionBox}
\end{frame}

% SLIDE ID: reg-ctms
\begin{frame}{CTMS}
  \SlideID{reg-ctms}
  \begin{itemize}
    \item CTMS wspiera nadzór nad badaniem i ośrodkami,
    \item statusy i KPI powinny być spójne z rzeczywistością,
    \item działania monitorujące muszą być dokumentowane,
    \item integracje z EDC i eTMF powinny być opisane i kontrolowane.
  \end{itemize}

  \vspace{0.3cm}

  \begin{ExampleBox}
    Jeśli CTMS pokazuje „site active”, system powinien wspierać dowód, że wymagane kroki aktywacji naprawdę wykonano.
  \end{ExampleBox}
\end{frame}

% SLIDE ID: reg-epro
\begin{frame}{ePRO / eCOA}
  \SlideID{reg-epro}
  \begin{itemize}
    \item dane pacjenta muszą być przypisane do właściwej osoby,
    \item czas wpisu i czas obserwacji mają znaczenie,
    \item aplikacja powinna ograniczać brakujące i opóźnione wpisy,
    \item zgłoszenie istotnego objawu może wymagać alertu do zespołu badania.
  \end{itemize}

  \vspace{0.3cm}

  \begin{WarningBox}
    Pacjent zgłasza duszność w dzienniczku. System nie powinien być tylko „notatnikiem”.
  \end{WarningBox}
\end{frame}

% SLIDE ID: reg-vendor-saas
\begin{frame}{Vendor i SaaS}
  \SlideID{reg-vendor-saas}
  \begin{itemize}
    \item odpowiedzialność nie znika po wyborze dostawcy,
    \item umowa powinna opisywać role, dane, dostęp i dokumentację,
    \item konfiguracja produkcyjna musi być znana i kontrolowana,
    \item zmiany, aktualizacje i incydenty wymagają procesu.
  \end{itemize}

  \vspace{0.3cm}

  \begin{WarningBox}
    „System jest w chmurze” nie oznacza automatycznie, że jest zgodny z GCP.
  \end{WarningBox}
\end{frame}

% SLIDE ID: reg-ai-note
\begin{frame}{AI w systemach badania}
  \SlideID{reg-ai-note}
  \begin{itemize}
    \item AI może wspierać rekrutację, coding, query i monitoring,
    \item wynik AI musi być możliwy do oceny przez człowieka,
    \item trzeba kontrolować bias, jakość danych i zmianę modelu,
    \item odpowiedzialność za decyzję pozostaje po stronie ludzi i organizacji.
  \end{itemize}

  \vspace{0.3cm}

  \begin{QuestionBox}
    Kiedy sugestia AI staje się decyzją kliniczną albo operacyjną?
  \end{QuestionBox}
\end{frame}

% SLIDE ID: reg-practical-checklist
\begin{frame}{Checklist dla systemu}
  \SlideID{reg-practical-checklist}
  \begin{itemize}
    \item Czy wiemy, jakie dane są źródłowe?
    \item Czy każdy użytkownik ma indywidualne konto?
    \item Czy audit trail jest dostępny i czytelny?
    \item Czy mamy backup, archiwizację i plan awaryjny?
    \item Czy system i konfiguracja są zwalidowane?
  \end{itemize}
\end{frame}

% SLIDE ID: reg-summary
\begin{frame}{Najważniejszy wniosek}
  \SlideID{reg-summary}
  \begin{itemize}
    \item regulacje nie wymagają „jednego idealnego systemu”,
    \item wymagają kontroli nad ryzykiem, danymi i procesem,
    \item funkcja w aplikacji musi być wsparta SOP, walidacją i szkoleniem,
    \item im większy wpływ na pacjenta i wyniki, tym większy poziom kontroli.
  \end{itemize}
\end{frame}

% SLIDE ID: reg-sources
\begin{frame}{Źródła}
  \SlideID{reg-sources}
  \footnotesize
  \begin{itemize}
    \item ICH E6(R3), \textit{Good Clinical Practice}, final guideline, 2025.
    \item EMA, \textit{Guideline on computerised systems and electronic data in clinical trials}, EMA/INS/GCP/112288/2023.
    \item FDA, \textit{Electronic Systems, Electronic Records, and Electronic Signatures in Clinical Investigations: Questions and Answers}, 2024.
    \item Regulation (EU) No 536/2014 on clinical trials on medicinal products for human use.
    \item EMA, \textit{ICH E6 Good Clinical Practice — Scientific guideline}.
  \end{itemize}
\end{frame}
